\documentclass{amsart}
% For dealing with references we use the comment environment
\usepackage{verbatim}
\newenvironment{reference}{\comment}{\endcomment}
%\newenvironment{reference}{}{}
\newenvironment{slogan}{\comment}{\endcomment}
\newenvironment{history}{\comment}{\endcomment}

% For commutative diagrams we use Xy-pic
\usepackage[all]{xy}

% We use 2cell for 2-commutative diagrams.
\xyoption{2cell}
\UseAllTwocells

% We use multicol for the list of chapters between chapters
\usepackage{multicol}

% This is generally recommended for better output
\usepackage{lmodern}
\usepackage[T1]{fontenc}

% For cross-file-references

% Package for hypertext links:
\usepackage{hyperref}

% For any local file, say "hello.tex" you want to link to please

% Theorem environments.
%
\theoremstyle{plain}
\newtheorem{theorem}[subsection]{Theorem}
\newtheorem{proposition}[subsection]{Proposition}
\newtheorem{lemma}[subsection]{Lemma}

\theoremstyle{definition}
\newtheorem{definition}[subsection]{Definition}
\newtheorem{example}[subsection]{Example}
\newtheorem{exercise}[subsection]{Exercise}
\newtheorem{situation}[subsection]{Situation}

\theoremstyle{remark}
\newtheorem{remark}[subsection]{Remark}
\newtheorem{remarks}[subsection]{Remarks}

\numberwithin{equation}{subsection}

% Macros
%
\def\lim{\mathop{\mathrm{lim}}\nolimits}
\def\colim{\mathop{\mathrm{colim}}\nolimits}
\def\Spec{\mathop{\mathrm{Spec}}}
\def\Hom{\mathop{\mathrm{Hom}}\nolimits}
\def\Ext{\mathop{\mathrm{Ext}}\nolimits}
\def\SheafHom{\mathop{\mathcal{H}\!\mathit{om}}\nolimits}
\def\SheafExt{\mathop{\mathcal{E}\!\mathit{xt}}\nolimits}
\def\Sch{\mathit{Sch}}
\def\Mor{\mathop{\mathrm{Mor}}\nolimits}
\def\Ob{\mathop{\mathrm{Ob}}\nolimits}
\def\Sh{\mathop{\mathit{Sh}}\nolimits}
\def\NL{\mathop{N\!L}\nolimits}
\def\CH{\mathop{\mathrm{CH}}\nolimits}
\def\proetale{{pro\text{-}\acute{e}tale}}
\def\etale{{\acute{e}tale}}
\def\QCoh{\mathit{QCoh}}
\def\Ker{\mathop{\mathrm{Ker}}}
\def\Im{\mathop{\mathrm{Im}}}
\def\Coker{\mathop{\mathrm{Coker}}}
\def\Coim{\mathop{\mathrm{Coim}}}

% Boxtimes
%
\DeclareMathSymbol{\boxtimes}{\mathbin}{AMSa}{"02}

%
% Macros for moduli stacks/spaces
%
\def\QCohstack{\mathcal{QC}\!\mathit{oh}}
\def\Cohstack{\mathcal{C}\!\mathit{oh}}
\def\Spacesstack{\mathcal{S}\!\mathit{paces}}
\def\Quotfunctor{\mathrm{Quot}}
\def\Hilbfunctor{\mathrm{Hilb}}
\def\Curvesstack{\mathcal{C}\!\mathit{urves}}
\def\Polarizedstack{\mathcal{P}\!\mathit{olarized}}
\def\Complexesstack{\mathcal{C}\!\mathit{omplexes}}
% \Pic is the operator that assigns to X its picard group, usage \Pic(X)
% \Picardstack_{X/B} denotes the Picard stack of X over B
% \Picardfunctor_{X/B} denotes the Picard functor of X over B
\def\Pic{\mathop{\mathrm{Pic}}\nolimits}
\def\Picardstack{\mathcal{P}\!\mathit{ic}}
\def\Picardfunctor{\mathrm{Pic}}
\def\Deformationcategory{\mathcal{D}\!\mathit{ef}}

\setlength{\emergencystretch}{3em}
\begin{document}
\title[Stacks Project, Tag 07F9]{732 --- Smooth local resolutions delocalize at primes minimal over the singular ideal}
\maketitle
\noindent Copyright (C) 2005--2025 Johan de Jong. Stacks Project authors. Source: \href{https://stacks.math.columbia.edu/tag/07F9}{Tag 07F9}.

\noindent Editorial difficulty note (not author text). Difficulty H2: The proof homogenizes a localized standard-smooth presentation, controls denominator annihilators, and constructs a global finite-presentation factorization with a smooth open avoiding the chosen prime. This explicit delocalization construction is beyond qualification commutative algebra..

\noindent Editorial provenance note (not author text). History: excerpt from revision \texttt{a0444\allowbreak 6e57e\allowbreak c1fbc\allowbreak 252a8\allowbreak 71afc\allowbreak ec775\allowbreak 2fb28\allowbreak 07b14}, smoothing.tex, lines 2147--2227. Reference presentation was adapted; original-author context and core support blocks were retained or added. The mathematical wording is unchanged. Licensed under GNU FDL 1.2 or later; no invariant sections or cover texts. See ../licenses/COPYING. Context blocks reproduce author definitions and supporting results; they are not additional counted problems.

\begin{situation}
\label{situation-local}
We are given a Noetherian ring $R$ and an $R$-algebra map $A \to \Lambda$
and a prime $\mathfrak q \subset \Lambda$. We assume $A$ is of
finite presentation over $R$. In this situation we denote
$\mathfrak h_A = \sqrt{H_{A/R} \Lambda}$.
\end{situation}

\begin{definition}
\label{definition-singular-ideal}
Let $R \to A$ be a ring map. The {\it singular ideal of $A$ over $R$},
denoted $H_{A/R}$ is the unique radical ideal $H_{A/R} \subset A$ with
$$
V(H_{A/R}) = \{\mathfrak q \in \Spec(A) \mid R \to A
\text{ not smooth at }\mathfrak q\}
$$
\end{definition}

\begin{definition}
\label{algebra-definition-unramified}
Let $R \to S$ be a ring map.
\begin{enumerate}
\item We say $R \to S$ is {\it unramified} if $R \to S$ is of
finite type and $\Omega_{S/R} = 0$.
\item We say $R \to S$ is {\it G-unramified} if $R \to S$ is of finite
presentation and $\Omega_{S/R} = 0$.
\item Given a prime $\mathfrak q$ of $S$ we say that $S$ is
{\it unramified at $\mathfrak q$} if there exists a
$g \in S$, $g \not \in \mathfrak q$ such that $R \to S_g$ is unramified.
\item Given a prime $\mathfrak q$ of $S$ we say that $S$ is
{\it G-unramified at $\mathfrak q$} if there exists a
$g \in S$, $g \not \in \mathfrak q$ such that $R \to S_g$ is G-unramified.
\end{enumerate}
\end{definition}

\begin{lemma}
\label{lemma-smooth-standard-smooth}
Let $R \to A$ be a smooth ring map. Then there exists a smooth $R$-algebra
map $A \to B$ with a retraction such that $B$ is standard smooth over
$R$, i.e.,
$$
B \cong R[x_1, \ldots, x_n]/(f_1, \ldots, f_c)
$$
and $\det(\partial f_j/\partial x_i)_{i, j = 1, \ldots, c}$
is invertible in $B$.
\end{lemma}

\begin{proof}
Apply Lemma \href{https://stacks.math.columbia.edu/tag/07CG}{lemma syntomic complete intersection (Tag 07CG)}
to get a smooth $R$-algebra map $A \to C$ with a retraction such that
$C = R[x_1, \ldots, x_n]/(f_1, \ldots, f_c)$
is a relative global complete intersection over $R$. As $C$ is smooth
over $R$ we have a short exact sequence
$$
0 \to
\bigoplus\nolimits_{j = 1, \ldots, c} C f_j \to
\bigoplus\nolimits_{i = 1, \ldots, n} C\text{d}x_i \to
\Omega_{C/R} \to 0
$$
Since $\Omega_{C/R}$ is a projective $C$-module this sequence is split.
Choose a left inverse $t$ to the first map. Say
$t(\text{d}x_i) = \sum c_{ij} f_j$
so that $\sum_i \frac{\partial f_j}{\partial x_i} c_{i\ell} = \delta_{j\ell}$
(Kronecker delta). Let
$$
B' = C[y_1, \ldots, y_c] =
R[x_1, \ldots, x_n, y_1, \ldots, y_c]/(f_1, \ldots, f_c)
$$
The $R$-algebra map $C \to B'$ has a retraction given by mapping $y_j$ to zero.
We claim that the map
$$
R[z_1, \ldots, z_n] \longrightarrow B',\quad
z_i \longmapsto x_i - \sum\nolimits_j c_{ij} y_j
$$
is \'etale at every point in the image of $\Spec(C) \to \Spec(B')$.
In $\Omega_{B'/R[z_1, \ldots, z_n]}$ we have
$$
0 =
\text{d}f_j - \sum\nolimits_i \frac{\partial f_j}{\partial x_i} \text{d}z_i
\equiv
\sum\nolimits_{i, \ell}
\frac{\partial f_j}{\partial x_i} c_{i\ell} \text{d}y_\ell
\equiv
\text{d}y_j \bmod (y_1, \ldots, y_c)\Omega_{B'/R[z_1, \ldots, z_n]}
$$
Since $0 = \text{d}z_i = \text{d}x_i$ modulo
$\sum B'\text{d}y_j + (y_1, \ldots, y_c)\Omega_{B'/R[z_1, \ldots, z_n]}$
we conclude that
$$
\Omega_{B'/R[z_1, \ldots, z_n]}/
(y_1, \ldots, y_c)\Omega_{B'/R[z_1, \ldots, z_n]} = 0.
$$
As $\Omega_{B'/R[z_1, \ldots, z_n]}$ is a finite $B'$-module
by Nakayama's lemma there exists a $g \in 1 + (y_1, \ldots, y_c)$
that $(\Omega_{B'/R[z_1, \ldots, z_n]})_g = 0$. This proves that
$R[z_1, \ldots, z_n] \to B'_g$ is unramified, see
Algebra, Definition \ref{algebra-definition-unramified}.
For any ring map $R \to k$ where $k$ is a field we obtain an
unramified ring map $k[z_1, \ldots, z_n] \to (B'_g) \otimes_R k$
between smooth $k$-algebras of dimension $n$. It follows that
$k[z_1, \ldots, z_n] \to (B'_g) \otimes_R k$ is flat by
Algebra, Lemmas \href{https://stacks.math.columbia.edu/tag/00R4}{lemma CM over regular flat (Tag 00R4)} and
\href{https://stacks.math.columbia.edu/tag/00TS}{lemma characterize smooth kbar (Tag 00TS)}. By the crit\`ere
de platitude par fibre
(Algebra, Lemma \href{https://stacks.math.columbia.edu/tag/00R7}{lemma criterion flatness fibre (Tag 00R7)})
we conclude that $R[z_1, \ldots, z_n] \to B'_g$ is flat.
Finally, Algebra, Lemma \href{https://stacks.math.columbia.edu/tag/00U6}{lemma characterize etale (Tag 00U6)}
implies that $R[z_1, \ldots, z_n] \to B'_g$ is \'etale.
Set $B = B'_g$. Note that $C \to B$ is smooth and has a retraction,
so also $A \to B$ is smooth and has a retraction.
Moreover, $R[z_1, \ldots, z_n] \to B$ is \'etale.
By Algebra, Lemma \href{https://stacks.math.columbia.edu/tag/00U9}{lemma etale standard smooth (Tag 00U9)}
we can write
$$
B = R[z_1, \ldots, z_n, w_1, \ldots, w_m]/(g_1, \ldots, g_m)
$$
with $\det(\partial g_j/\partial w_i)$ invertible in $B$.
This proves the lemma.
\end{proof}


Let $R \to A \to \Lambda \supset \mathfrak q$ be as in
Situation \ref{situation-local}. We say
{\it $R \to A \to \Lambda \supset \mathfrak q$
can be resolved} if there exists a factorization $A \to B \to \Lambda$
with $B$ of finite presentation and
$\mathfrak h_A \subset \mathfrak h_B \not \subset \mathfrak q$.
In this case we will call the factorization $A \to B \to \Lambda$
a {\it resolution of $R \to A \to \Lambda \supset \mathfrak q$}.


\begin{lemma}
\label{lemma-delocalize-weak}
\begin{reference}
\href{https://stacks.math.columbia.edu/bibliography}{Bibliography: swan (Lemma 12.2)} or
\href{https://stacks.math.columbia.edu/bibliography}{Bibliography: popescu-GND (Lemma 2)}
\end{reference}
Let $R \to A \to \Lambda \supset \mathfrak q$ be as in
Situation \ref{situation-local}. Let $\mathfrak p = R \cap \mathfrak q$.
Assume that $\mathfrak q$ is minimal over $\mathfrak h_A$ and that
$R_\mathfrak p \to A_\mathfrak p \to \Lambda_\mathfrak q
\supset \mathfrak q\Lambda_\mathfrak q$ can be resolved.
Then there exists a factorization $A \to C \to \Lambda$ with $C$ of
finite presentation such that $H_{C/R} \Lambda \not \subset \mathfrak q$.
\end{lemma}

\begin{proof}
Let $A_\mathfrak p \to C \to \Lambda_\mathfrak q$ be a resolution of
$R_\mathfrak p \to A_\mathfrak p \to \Lambda_\mathfrak q
\supset \mathfrak q\Lambda_\mathfrak q$. By our assumption
that $\mathfrak q$ is minimal over $\mathfrak h_A$ this
means that $H_{C/R_\mathfrak p} \Lambda_\mathfrak q = \Lambda_\mathfrak q$.
By Lemma \href{https://stacks.math.columbia.edu/tag/07EU}{lemma final solve (Tag 07EU)}
we may assume that $C$ is smooth over $R_\mathfrak p$.
By Lemma \ref{lemma-smooth-standard-smooth} we may assume that
$C$ is standard smooth over $R_\mathfrak p$.
Write $A = R[x_1, \ldots, x_n]/(g_1, \ldots, g_t)$ and say
$A \to \Lambda$ is given by $x_i \mapsto \lambda_i$.
Write $C = R_\mathfrak p[x_1, \ldots, x_{n + m}]/(f_1, \ldots, f_c)$
for some $c \geq n$ such that $A \to C$ maps $x_i$ to $x_i$ and such that
$\det(\partial f_j/\partial x_i)_{i, j = 1, \ldots, c}$
is invertible in $C$, see
Lemma \href{https://stacks.math.columbia.edu/tag/07EY}{lemma standard smooth include generators (Tag 07EY)}.
After clearing denominators we may assume
$f_1, \ldots, f_c$ are elements of $R[x_1, \ldots, x_{n + m}]$.
Of course
$\det(\partial f_j/\partial x_i)_{i, j = 1, \ldots, c}$
is not invertible in $R[x_1, \ldots, x_{n + m}]/(f_1, \ldots, f_c)$
but it becomes invertible after inverting some element $s_0 \in R$,
$s_0 \not \in \mathfrak p$.
As $g_j$ maps to zero under $R[x_1, \ldots, x_n] \to A \to C$
we can find $s_j \in R$, $s_j \not \in \mathfrak p$ such that
$s_j g_j$ is zero in $R[x_1, \ldots, x_{n + m}]/(f_1, \ldots, f_c)$.
Write $f_j = F_j(x_1, \ldots, x_{n + m}, 1)$
for some polynomial
$F_j \in R[x_1, \ldots, x_n, X_{n + 1}, \ldots, X_{n + m + 1}]$
homogeneous in $X_{n + 1}, \ldots, X_{n + m + 1}$.
Pick $\lambda_{n + i} \in \Lambda$, $i = 1, \ldots, m + 1$ with
$\lambda_{n + m + 1} \not \in \mathfrak q$ such that $x_{n + i}$ maps to
$\lambda_{n + i}/\lambda_{n + m + 1}$ in $\Lambda_\mathfrak q$.
Then
\begin{align*}
F_j(\lambda_1, \ldots, \lambda_{n + m + 1})
& =
(\lambda_{n + m + 1})^{\deg(F_j)} F_j(\lambda_1, \ldots, \lambda_n,
\frac{\lambda_{n + 1}}{\lambda_{n + m + 1}}, \ldots,
\frac{\lambda_{n + m}}{\lambda_{n + m + 1}}, 1) \\
& =
(\lambda_{n + m + 1})^{\deg(F_j)} f_j(\lambda_1, \ldots, \lambda_n,
\frac{\lambda_{n + 1}}{\lambda_{n + m + 1}}, \ldots,
\frac{\lambda_{n + m}}{\lambda_{n + m + 1}}) \\
& = 0
\end{align*}
in $\Lambda_\mathfrak q$. Thus we can find
$\lambda_0 \in \Lambda$, $\lambda_0 \not \in \mathfrak q$ such that
$\lambda_0 F_j(\lambda_1, \ldots, \lambda_{n + m + 1}) = 0$
in $\Lambda$. Now we set $B$ equal to
$$
R[x_0, \ldots, x_{n + m + 1}]/
(g_1, \ldots, g_t, x_0F_1(x_1, \ldots, x_{n + m + 1}), \ldots,
x_0F_c(x_1, \ldots, x_{n + m + 1}))
$$
which we map to $\Lambda$ by mapping $x_i$ to $\lambda_i$.
Let $b$ be the image of $x_0 x_{n + m + 1} s_0 s_1 \ldots s_t$ in $B$.
Then $B_b$ is isomorphic to
$$
R_{s_0s_1 \ldots s_t}[x_0, x_1, \ldots, x_{n + m + 1}, 1/x_0x_{n + m + 1}]/
(f_1, \ldots, f_c)
$$
which is smooth over $R$ by construction.
Since $b$ does not map to an element of $\mathfrak q$, we win.
\end{proof}
\end{document}
